\uchapter{Abstract}

In an increasingly competitive logistics landscape, the ability to turn raw operational data
into actionable decision-making insight has become a strategic imperative for logistics service
providers. Yet many organisations still rely on fragmented information systems that lack
integration, real-time visibility, and analytical depth --- leaving decision-makers without the
means to monitor performance and control costs effectively.

This thesis presents the design and deployment of a Business Intelligence (BI) decision-support
system for the analysis of logistics operations and costs within a multi-site logistics network.
The work is built around a centralised Data Warehouse (DW) as its analytical backbone, fed by an
automated Extract-Transform-Load (ETL) pipeline that consolidates data from heterogeneous source
systems, and completed by interactive dashboards and a proactive alerting mechanism that support
timely, well-informed decisions.

The thesis treats two business axes: \emph{on-demand dedicated transport}, covering the analysis
of expressed demand and the evaluation of associated cost and profitability; and \emph{parcel
delivery}, covering the e-commerce express parcel lifecycle at scale together with its logistics
cost and profitability. For each business axis, three aspects are analysed: operations, cost and
profitability, and performance.

Overall, this thesis bridges the gap between raw logistics data and strategic decision support,
providing logistics managers with a reliable foundation for operational excellence and cost
governance.

\bigskip
\vspace{0.5cm}

\noindent\textbf{Keywords:} Business Intelligence, Data Warehouse, ETL, Dashboard, Logistics.
